\documentclass[10pt,a4paper]{amsart}
\usepackage{amsmath,amssymb,amsthm,graphicx}
\usepackage{verbatim}
\newenvironment{reference}{\comment}{\endcomment}
%\newenvironment{reference}{}{}
\newenvironment{slogan}{\comment}{\endcomment}
\newenvironment{history}{\comment}{\endcomment}
% For commutative diagrams we use Xy-pic
\usepackage[all]{xy}

% We use 2cell for 2-commutative diagrams.
\xyoption{2cell}
% Theorem environments.
%
\theoremstyle{plain}
\newtheorem{theorem}[subsection]{Theorem}
\newtheorem{proposition}[subsection]{Proposition}
\newtheorem{lemma}[subsection]{Lemma}

\theoremstyle{definition}
\newtheorem{definition}[subsection]{Definition}
\newtheorem{example}[subsection]{Example}
\newtheorem{exercise}[subsection]{Exercise}
\newtheorem{situation}[subsection]{Situation}

\theoremstyle{remark}
\newtheorem{remark}[subsection]{Remark}
\newtheorem{remarks}[subsection]{Remarks}

\numberwithin{equation}{subsection}

% Macros
%
\def\lim{\mathop{\mathrm{lim}}\nolimits}
\def\colim{\mathop{\mathrm{colim}}\nolimits}
\def\Spec{\mathop{\mathrm{Spec}}}
\def\Hom{\mathop{\mathrm{Hom}}\nolimits}
\def\Ext{\mathop{\mathrm{Ext}}\nolimits}
\def\SheafHom{\mathop{\mathcal{H}\!\mathit{om}}\nolimits}
\def\SheafExt{\mathop{\mathcal{E}\!\mathit{xt}}\nolimits}
\def\Sch{\mathit{Sch}}
\def\Mor{\mathop{\mathrm{Mor}}\nolimits}
\def\Ob{\mathop{\mathrm{Ob}}\nolimits}
\def\Sh{\mathop{\mathit{Sh}}\nolimits}
\def\NL{\mathop{N\!L}\nolimits}
\def\CH{\mathop{\mathrm{CH}}\nolimits}
\def\proetale{{pro\text{-}\acute{e}tale}}
\def\etale{{\acute{e}tale}}
\def\QCoh{\mathit{QCoh}}
\def\Ker{\mathop{\mathrm{Ker}}}
\def\Im{\mathop{\mathrm{Im}}}
\def\Coker{\mathop{\mathrm{Coker}}}
\def\Coim{\mathop{\mathrm{Coim}}}

% Boxtimes
%
\DeclareMathSymbol{\boxtimes}{\mathbin}{AMSa}{"02}

%
% Macros for moduli stacks/spaces
%
\def\QCohstack{\mathcal{QC}\!\mathit{oh}}
\def\Cohstack{\mathcal{C}\!\mathit{oh}}
\def\Spacesstack{\mathcal{S}\!\mathit{paces}}
\def\Quotfunctor{\mathrm{Quot}}
\def\Hilbfunctor{\mathrm{Hilb}}
\def\Curvesstack{\mathcal{C}\!\mathit{urves}}
\def\Polarizedstack{\mathcal{P}\!\mathit{olarized}}
\def\Complexesstack{\mathcal{C}\!\mathit{omplexes}}
% \Pic is the operator that assigns to X its picard group, usage \Pic(X)
% \Picardstack_{X/B} denotes the Picard stack of X over B
% \Picardfunctor_{X/B} denotes the Picard functor of X over B
\def\Pic{\mathop{\mathrm{Pic}}\nolimits}
\def\Picardstack{\mathcal{P}\!\mathit{ic}}
\def\Picardfunctor{\mathrm{Pic}}
\def\Deformationcategory{\mathcal{D}\!\mathit{ef}}
\usepackage{cleveref}
\crefname{originallemma}{original lemma}{original lemmas}
\Crefname{originallemma}{Original lemma}{Original lemmas}
\makeatletter
\newlabel{lemma-abelian}{{lemma-abelian}{source}{}{sourcepointer-1}{}}
\newlabel{lemma-abelian@cref}{{[originallemma][0][]lemma-abelian}{[1][0][]source}{}{}{}}
\makeatother
\crefname{originallemma}{original lemma}{original lemmas}
\Crefname{originallemma}{Original lemma}{Original lemmas}
\makeatletter
\newlabel{lemma-stalk-module-differentials}{{lemma-stalk-module-differentials}{source}{}{sourcepointer-2}{}}
\newlabel{lemma-stalk-module-differentials@cref}{{[originallemma][0][]lemma-stalk-module-differentials}{[1][0][]source}{}{}{}}
\makeatother
\crefname{originalreference}{original reference}{original references}
\Crefname{originalreference}{Original reference}{Original references}
\makeatletter
\newlabel{algebra-lemma-differential-seq}{{algebra-lemma-differential-seq}{source}{}{sourcepointer-3}{}}
\newlabel{algebra-lemma-differential-seq@cref}{{[originalreference][0][]algebra-lemma-differential-seq}{[1][0][]source}{}{}{}}
\makeatother
\title{The exact sequence relating the conormal sheaf of a surjection of sheaves of algebras to the modules of differentials}
\author{The Stacks Project Authors}
\date{stacks-project:modules.tex:local-snapshot-20261012; GNU FDL 1.2 or later}
\begin{document}
\maketitle
\noindent\textit{Editorial note.} All delivered mathematical text is the author's own, copied from the pinned Stacks Project chapter source sources/stacks-modules/pinned-original-source.tex; no statement or proof step was rewritten, reordered or removed.\par\medskip
\noindent\textit{Editorial note.} Preamble provenance: the Stacks Project keeps one single shared preamble for every chapter; this card takes it directly from the staged sources/stacks-modules/preamble.tex as three verbatim spans (lines 8-12, 14-18 and 154-217), so no derived file was written and nothing was added or redefined. The wrapper class and the base packages are supplied by the assembler.\par\medskip
\noindent\textit{Editorial note.} Bibliography provenance: neither the counted statement nor its proof contains a citation (this lane rescans the whole span before writing the card), so the card carries no bibliography and no entry was reconstructed.\par\medskip
\noindent\textit{Editorial note.} Selection provenance: this card was opened under the file-based control contract (CONTROL.md), with the mechanical pre-card gate run on its statement and proof interval before the card was written.\par\medskip
\noindent\textit{Allowed standard prerequisites (editorial).}\begin{itemize}
\item The functoriality of the module of differentials for a commutative diagram of sheaves of rings, with the map characterised on differentials of local sections. Author lemma lemma-functoriality-differentials, whose statement is reproduced verbatim as a context span of this card (source lines 5136-5157). Supplies the situation in which the statement is formulated and the right hand map of the sequence.
\item The Leibniz rule for the universal derivation, and the compatibility of the formation of differentials with taking stalks. Used inside the counted proof (source lines 5184-5191 and 5195-5197). Shows that the displayed rule defines a map on the quotient by the square of the kernel and identifies the sequence with the algebraic one.
\end{itemize}\medskip
\ifdefined\hypertarget\hypertarget{sourcepointer-1}{}\fi
\noindent\textit{Source reference.} Original lemma lemma-abelian: sources/stacks-modules/pinned-original-source.tex, source lines 439--439. Navigation pointer only: the counted text refers to the Stacks item lemma-abelian, which belongs to another chapter of the project and is not part of this card. The pointer carries the target's own label as its number because this chapter's staged snapshot does not contain the other chapter's tag table; its content is not reproduced here.\par\medskip
\ifdefined\hypertarget\hypertarget{sourcepointer-2}{}\fi
\noindent\textit{Source reference.} Original lemma lemma-stalk-module-differentials: sources/stacks-modules/pinned-original-source.tex, source lines 5197--5197. Navigation pointer only: the counted text refers to the Stacks item lemma-stalk-module-differentials, which belongs to another chapter of the project and is not part of this card. The pointer carries the target's own label as its number because this chapter's staged snapshot does not contain the other chapter's tag table; its content is not reproduced here.\par\medskip
\ifdefined\hypertarget\hypertarget{sourcepointer-3}{}\fi
\noindent\textit{Source reference.} Original reference algebra-lemma-differential-seq: sources/stacks-modules/pinned-original-source.tex, source lines 5199--5199. Navigation pointer only: the counted text refers to the Stacks item algebra-lemma-differential-seq, which belongs to another chapter of the project and is not part of this card. The pointer carries the target's own label as its number because this chapter's staged snapshot does not contain the other chapter's tag table; its content is not reproduced here.\par\medskip
\begin{lemma}
\label{lemma-differential-seq}
In Lemma \ref{lemma-functoriality-differentials} suppose that
$\mathcal{O}_2 \to \mathcal{O}'_2$ is surjective with kernel
$\mathcal{I} \subset \mathcal{O}_2$ and assume that
$\mathcal{O}_1 = \mathcal{O}'_1$. Then there is a canonical exact
sequence of $\mathcal{O}'_2$-modules
$$
\mathcal{I}/\mathcal{I}^2
\longrightarrow
\Omega_{\mathcal{O}_2/\mathcal{O}_1} \otimes_{\mathcal{O}_2} \mathcal{O}'_2
\longrightarrow
\Omega_{\mathcal{O}'_2/\mathcal{O}_1}
\longrightarrow
0
$$
The leftmost map is characterized by the rule that a local section
$f$ of $\mathcal{I}$ maps to $\text{d}f \otimes 1$.
\end{lemma}

\begin{lemma}
\label{lemma-functoriality-differentials}
Let $X$ be a topological space. Let
$$
\xymatrix{
\mathcal{O}_2 \ar[r]_\varphi & \mathcal{O}_2' \\
\mathcal{O}_1 \ar[r] \ar[u] & \mathcal{O}'_1 \ar[u]
}
$$
be a commutative diagram of sheaves of rings on $X$. The map
$\mathcal{O}_2 \to \mathcal{O}'_2$ composed with the map
$\text{d} : \mathcal{O}'_2 \to \Omega_{\mathcal{O}'_2/\mathcal{O}'_1}$
is a $\mathcal{O}_1$-derivation. Hence we obtain a canonical map of
$\mathcal{O}_2$-modules
$\Omega_{\mathcal{O}_2/\mathcal{O}_1} \to
\Omega_{\mathcal{O}'_2/\mathcal{O}'_1}$.
It is uniquely characterized by the property that
$\text{d}(f) \mapsto \text{d}(\varphi(f))$
for any local section $f$ of $\mathcal{O}_2$.
In this way $\Omega_{-/-}$ becomes a functor on the category
of arrows of sheaves of rings.
\end{lemma}

\begin{proof}
For a local section $f$ of $\mathcal{I}$ denote $\overline{f}$ the image of
$f$ in $\mathcal{I}/\mathcal{I}^2$. To show that the map
$\overline{f} \mapsto \text{d}f \otimes 1$ is well defined we just have to
check that $\text{d} f_1f_2 \otimes 1 = 0$ if $f_1, f_2$ are local sections
of $\mathcal{I}$. And this is clear from the Leibniz rule
$\text{d} f_1f_2 \otimes 1 =
(f_1 \text{d}f_2 + f_2 \text{d} f_1 )\otimes 1 =
\text{d}f_2 \otimes f_1 + \text{d}f_1 \otimes f_2 = 0$.
A similar computation show this map is
$\mathcal{O}'_2 = \mathcal{O}_2/\mathcal{I}$-linear. The map on the right
is the one from Lemma \ref{lemma-functoriality-differentials}. To see
that the sequence is exact, we can check on stalks
(Lemma \ref{lemma-abelian}). By
Lemma \ref{lemma-stalk-module-differentials}
this follows from
Algebra, Lemma \ref{algebra-lemma-differential-seq}.
\end{proof}

\end{document}
